\textbf{结论名称：} 双曲线外一点引两切线的切点弦方程。

\textbf{适用条件：} 双曲线为标准方程 \(\frac{x^2}{a^2}-\frac{y^2}{b^2}=1\) 或 \(\frac{y^2}{a^2}-\frac{x^2}{b^2}=1\)，点 \(P(x_0,y_0)\) 在双曲线外部（可引两条实切线），两切点确定的直线称为切点弦。

\textbf{变量说明：} \(a\) 为实半轴长，\(b\) 为虚半轴长；\((x_0,y_0)\) 是双曲线外一点坐标；\((x,y)\) 为切点弦上动点坐标。

\textbf{核心公式：} 切点弦所在直线方程为
\[
\boxed{\frac{x_0 x}{a^2} - \frac{y_0 y}{b^2} = 1}
\quad\text{或}\quad
\boxed{\frac{y_0 y}{a^2} - \frac{x_0 x}{b^2} = 1}
\]
分别对应焦点在 \(x\) 轴与 \(y\) 轴的情形，与切线方程具有相同的代换规则。

\textbf{使用场景：} 已知双曲线和曲线外一点，快速写出切点弦方程而不需求解切点坐标；在极点极线、定点弦问题中可直接应用。

\textbf{简单例子：} 给定双曲线 \(\frac{x^2}{4}-\frac{y^2}{3}=1\) 及外一点 \(P(4,2)\)，代入第一式得切点弦方程：
\[
\frac{4x}{4} - \frac{2y}{3} = 1 \;\Longrightarrow\; x - \frac{2}{3}y = 1 \;\Longrightarrow\; 3x - 2y = 3.
\]